%!TEX TS-program = xelatex
%!TEX encoding = UTF-8 Unicode
% Awesome CV Hugo Bollon - Version d'une page optimisée
% Template: https://www.overleaf.com/latex/templates/awesome-cv/dfnvtnhzhhbm

\documentclass[11pt, a4paper]{awesome-cv}

% Configuration ------------------------------------------------
\geometry{left=1.2cm, top=1.0cm, right=1.2cm, bottom=1.2cm, footskip=.5cm}  % Marges réduites
\colorlet{awesome}{awesome-red}  % Choisir la couleur du thème
\setbool{acvSectionColorHighlight}{true}
\renewcommand{\acvHeaderSocialSep}{\quad\textbar\quad}  % Séparateur
\renewcommand{\acvSectionTopSkip}{4pt}
\renewcommand{\acvSectionContentTopSkip}{2mm}


% Données Personnelles -------------------------------------------------
\name{Hugo}{Bollon}
\position{Ingénieur Backend / DevOps \& Passionné d'Open Source}
\address{Grenoble, France}
\mobile{(+33) 6 51 99 80 78}
\email{hugo.bollon@proton.me}
\homepage{hugobollon.dev}
\github{hbollon}
\linkedin{hugobollon}
\extrainfo{Permis de conduire : B}

% Début du Document -----------------------------------------------
\begin{document}

%-------------------- Profil ----------------------------------
\makecvheader

\cvsection{Résumé}

\cvparagraph{
Ingénieur logiciel spécialisé en \textbf{Go} et passionné par l'\textbf{Open Source}, avec une solide expérience en \textbf{développement backend} et \textbf{ingénierie DevOps} (Kubernetes, AWS, CI/CD). Fort de plusieurs années d'expérience à concevoir et maintenir des plateformes cloud, je souhaite désormais mettre mes compétences et ma motivation au service de projets à fort impact en tant qu'\textbf{ingénieur Backend ou Fullstack}.
}

%-------------------- Compétences -----------------------------------
\cvsection{Compétences}

\begin{cvskills}
  \cvskill
    {Langages}
    {Go, Python, Java, C/C++, JavaScript, TypeScript, Bash}
  \cvskill
    {Frontend}
    {Vue, Angular}
  \cvskill
    {Backend}
    {API REST, Gorm, Gin, Echo, Gorilla Mux}
  \cvskill
    {DevOps}
    {Kubernetes, Docker, Helm, ArgoCD, Pulumi, Terraform, GitOps, AWS (EKS, RDS, EC2, ECS, Fargate), IaC, CI/CD}
  \cvskill
    {Base de données}
    {PostgreSQL, SQLite, Neo4j}
  \cvskill
    {IA / LLM Systems}
    {Opencode, Workflows agentiques, Architectures multi-agents, MCP, Local LLMs (llama.cpp, llama-swap)}
  \cvskill
    {Environnements}
    {Linux, Git, GitHub Actions, GitLab CI}
  \cvskill
    {Langues}
    {Français (langue maternelle), Anglais B2 (TOEIC 905), Japonais A1 (en apprentissage)}
\end{cvskills}

%-------------------- Expérience -------------------------------
\cvsection{Expérience}

\begin{cventries}
  \cventry
    {Ingénieur DevOps}
    {Camptocamp SA}
    {Chambéry, France}
    {Sep. 2021 - Présent}
    {
      \vspace{2.0mm}
      \begin{cvitems}
        \item {Conception, automatisation et maintenance de clusters Kubernetes avec plusieurs providers cloud (AWS, Exoscale, Scaleway) en utilisant l'Infrastructure as Code (Pulumi, Terraform) pour des plateformes fiables et évolutives.}
        \item {Développement backend en Go pour des outils internes, APIs et automatisations. Contributions open-source externes (providers Terraform, Zitadel, etc.).}
        \item {Référent open-source du département Infrastructure : gestion, maintenance et évolution des projets sous licence publique de Camptocamp Infrastructure Solutions.}
        \item {Conception et gestion de pipelines CI pour automatiser les étapes de build, lint, test et release, en utilisant Github Actions et Dagger.}
        \item {Mise en place de pipelines CD basés sur des pratiques GitOps avec ArgoCD, assurant des déploiements Kubernetes automatisés, reproductibles et sécurisés.}
        \item {Conteneurisation et orchestration de workloads : conception d'images Docker, configuration Docker Compose, déploiements orchestrés avec Kubernetes, et CI conteneurisée avec Dagger.}
        \item {Collaboration étroite avec le département Géospatial de Camptocamp : développement en Python et JavaScript (Vue 3, Angular), conception de charts Helm, et optimisation d'images Docker pour le déploiement de leurs applications.}
        \item {Déploiement et configuration de solutions de Monitoring, Logging et Alerting (Prometheus, Elastic, Grafana).}
        \item {Authentification OIDC fédérée (Keycloak, Zitadel).}
      \end{cvitems}
    }

  \cventry
    {Stagiaire Développeur Go}
    {Camptocamp SA}
    {Chambéry, France}
    {Juin 2021 - Sep. 2021}
    {
      \vspace{2.0mm}
      \begin{cvitems}
        \item {Développement et maintenance de Terraboard (outil d'observabilité Terraform).}
        \item {Migration du frontend d'AngularJS vers Vue 3 / TypeScript.}
        \item {Développement d'API backend en Go (Gorm, Gorilla Mux).}
        \item {Rédaction de tests unitaires et automatisation CI/CD.}
      \end{cvitems}
    }

  \cventry
    {Stagiaire Développeur Android}
    {LAET - ENTPE}
    {Chambéry, France}
    {Jan. 2020 - Mai 2020}
    {
      \vspace{2.0mm}
      \begin{cvitems}
        \item {Application Android native (Java) pour la digitalisation du jeu de plateau Urbalog avec support multi-devices à l'aide de l'API Nearby Connections.}
      \end{cvitems}
    }

\end{cventries}

%-------------------- Formation --------------------------------
\cvsection{Formation}

\begin{cventries}
  \cventry
    {Master Informatique (avec mention)}
    {Université Savoie Mont-Blanc}
    {Le-Bourget-Du-Lac, France}
    {2020 -- 2022}
    {}

  \cventry
    {Licence Informatique (avec mention)}
    {Université Savoie Mont-Blanc}
    {Le-Bourget-Du-Lac, France}
    {2017 -- 2020}
    {}

  \cventry
    {Baccalauréat Scientifique (ISN)}
    {Lycée de l'Albanais}
    {Rumilly, France}
    {2014 -- 2017}
    {}

\end{cventries}

%-------------------- Certifications ---------------------------
\cvsection{Certifications}

\begin{itemize}
  \item {\textbf{TOEIC} : 905/990 (Fév 2022) — Listening : 440/495 (B2), Reading : 465/495 (C1)}
  \item {\textbf{AWS Certified Cloud Practitioner} (Juin 2023) — ID : PST3Q8DCVE41Q5GD}
\end{itemize}

%-------------------- Atouts --------------------------------
\cvsection{Atouts}

\begin{tabular*}{\textwidth}{@{\extracolsep{\fill}}lcccccc}
Prise d'initiative & Autodidacte & Sociable & Attentif aux exigences & Apprentissage rapide & Esprit d'équipe \\
\end{tabular*}

\end{document}
